% GENERATED FILE. Edit canonical lessons, then run npm run generate:books.
% canonical-source-hash: fnv1a64:7047f964bee04318
% canonical-lessons: SA-C195-pattanam, SA-C195-rajadhani, SA-C195-janapadah, SA-C195-prantah, SA-C195-sima

\chapter{Port town, Capital, District, Province, Boundary}
\label{ch:sa-port-town-capital-district-province-boundary}
\hlchaptermodality{195}

\begin{chapteropening}
\textbf{By the end of this chapter:} \emph{I can say a port town, a capital, a district, a province and a boundary.}

Complete the last lesson of chapter 195: I can say a port town, a capital, a district, a province and a boundary.
\end{chapteropening}

\section[\texorpdfstring{\sk{पत्तनम्} (pattanam)}{पत्तनम् (pattanam)}]{\sk{पत्तनम्} (pattanam) — a port town}

\label{lesson:SA-C195-pattanam}



\begin{quote}
Before the new one: say the Sanskrit for a discussion group, then the Sanskrit for a fair.
\end{quote}

\subsection*{You'll want to know: \sk{पत्तनम्}}
\textbf{\sk{पत्तनम्}} — \emph{pattanam} — \textquotedblleft{}a port town\textquotedblright{}.

\begin{grammarlens}[title={What you've built}]
One more word for this chapter.
\end{grammarlens}

\subsection*{Your turn}
\begin{itemize}
\raggedright
  \item \emph{Say it:} \emph{pattanam}
  \item \emph{Say it:} \emph{pattanam}, once more
  \item \emph{Say it:} say \emph{pattanam}
  \item \emph{Recall:} say the Sanskrit for a discussion group, then the Sanskrit for a fair, then say \emph{pattanam} again
\end{itemize}

\begin{tcolorbox}[breakable,colback=teal!4,colframe=teal!35!black,title={Before you move on}]
What does \sk{पत्तनम्} mean? (\textquotedblleft{}A port town\textquotedblright{}.) Say it once more.
\end{tcolorbox}
\section[\texorpdfstring{\sk{राजधानी} (rājadhānī)}{राजधानी (rājadhānī)}]{\sk{राजधानी} (rājadhānī) — a capital}

\label{lesson:SA-C195-rajadhani}



\begin{quote}
Before the new one: say the Sanskrit for a fair, then the Sanskrit for a port town.
\end{quote}

\subsection*{You'll want to know: \sk{राजधानी}}
\textbf{\sk{राजधानी}} — \emph{rājadhānī} — \textquotedblleft{}a capital\textquotedblright{}.

\begin{grammarlens}[title={What you've built}]
One more word for this chapter.
\end{grammarlens}

\subsection*{Your turn}
\begin{itemize}
\raggedright
  \item \emph{Say it:} \emph{rājadhānī}
  \item \emph{Say it:} \emph{rājadhānī}, once more
  \item \emph{Say it:} say \emph{rājadhānī}
  \item \emph{Recall:} say the Sanskrit for a fair, then the Sanskrit for a port town, then say \emph{rājadhānī} again
\end{itemize}

\begin{tcolorbox}[breakable,colback=teal!4,colframe=teal!35!black,title={Before you move on}]
What does \sk{राजधानी} mean? (\textquotedblleft{}A capital\textquotedblright{}.) Say it once more.
\end{tcolorbox}
\section[\texorpdfstring{\sk{जनपदः} (janapadaḥ)}{जनपदः (janapadaḥ)}]{\sk{जनपदः} (janapadaḥ) — a district, a region}

\label{lesson:SA-C195-janapadah}



\begin{quote}
Before the new one: say the Sanskrit for a port town, then the Sanskrit for a capital.
\end{quote}

\subsection*{You'll want to know: \sk{जनपदः}}
\textbf{\sk{जनपदः}} — \emph{janapadaḥ} — \textquotedblleft{}a district, a region\textquotedblright{}.

\begin{grammarlens}[title={What you've built}]
One more word for this chapter.
\end{grammarlens}

\subsection*{Your turn}
\begin{itemize}
\raggedright
  \item \emph{Say it:} \emph{janapadaḥ}
  \item \emph{Say it:} \emph{janapadaḥ}, once more
  \item \emph{Say it:} say \emph{janapadaḥ}
  \item \emph{Recall:} say the Sanskrit for a port town, then the Sanskrit for a capital, then say \emph{janapadaḥ} again
\end{itemize}

\begin{tcolorbox}[breakable,colback=teal!4,colframe=teal!35!black,title={Before you move on}]
What does \sk{जनपदः} mean? (\textquotedblleft{}A district, a region\textquotedblright{}.) Say it once more.
\end{tcolorbox}
\section[\texorpdfstring{\sk{प्रान्तः} (prāntaḥ)}{प्रान्तः (prāntaḥ)}]{\sk{प्रान्तः} (prāntaḥ) — a province, a border}

\label{lesson:SA-C195-prantah}



\begin{quote}
Before the new one: say the Sanskrit for a capital, then the Sanskrit for a district.
\end{quote}

\subsection*{You'll want to know: \sk{प्रान्तः}}
\textbf{\sk{प्रान्तः}} — \emph{prāntaḥ} — \textquotedblleft{}a province, a border\textquotedblright{}.

\begin{grammarlens}[title={What you've built}]
One more word for this chapter.
\end{grammarlens}

\subsection*{Your turn}
\begin{itemize}
\raggedright
  \item \emph{Say it:} \emph{prāntaḥ}
  \item \emph{Say it:} \emph{prāntaḥ}, once more
  \item \emph{Say it:} say \emph{prāntaḥ}
  \item \emph{Recall:} say the Sanskrit for a capital, then the Sanskrit for a district, then say \emph{prāntaḥ} again
\end{itemize}

\begin{tcolorbox}[breakable,colback=teal!4,colframe=teal!35!black,title={Before you move on}]
What does \sk{प्रान्तः} mean? (\textquotedblleft{}A province, a border\textquotedblright{}.) Say it once more.
\end{tcolorbox}
\section[\texorpdfstring{\sk{सीमा} (sīmā)}{सीमा (sīmā)}]{\sk{सीमा} (sīmā) — a boundary}

\label{lesson:SA-C195-sima}



\begin{quote}
Before the new one: say the Sanskrit for a district, then the Sanskrit for a province.
\end{quote}

\subsection*{You'll want to know: \sk{सीमा}}
\textbf{\sk{सीमा}} — \emph{sīmā} — \textquotedblleft{}a boundary\textquotedblright{}.

\begin{grammarlens}[title={What you've built}]
One more word for this chapter.
\end{grammarlens}

\subsection*{Your turn}
\begin{itemize}
\raggedright
  \item \emph{Say it:} \emph{sīmā}
  \item \emph{Say it:} \emph{sīmā}, once more
  \item \emph{Say it:} say \emph{sīmā}
  \item \emph{Recall:} say the Sanskrit for a district, then the Sanskrit for a province, then say \emph{sīmā} again
\end{itemize}

\begin{tcolorbox}[breakable,colback=teal!4,colframe=teal!35!black,title={Before you move on}]
What does \sk{सीमा} mean? (\textquotedblleft{}A boundary\textquotedblright{}.) Say it once more.
\end{tcolorbox}
